\subsection{Brazilian Post-salt Field}

Results for the standard DSI-ESMDA are presented in \cref{tab:real-dsi-esmda},
while results for the proposed Gated DSI-ESMDA are shown in 
\cref{tab:real-gated-dsi-esmda-1-16}.

Unlike the trend set by the synthetic case studies,
localization did not have a significant impact on observation coverages or data mismatches in the real field case study,
be it in the standard DSI-ESMDA or the proposed Gated DSI-ESMDA.
%
However, it is clear that the gating produces better coverages as the gating intensity increases,
with the best results occuring at $\lambda=16$, the largest value tested.
Importantly, this comes at the cost of a higher data mismatch and MSE, which is expected since the posterior is gaining variance.

\begin{table}[htbp]
    \centering
    \IBGEtab{%
        \caption{%
            DSI-ESMDA results for the real field in the Brazilian Post-salt.
        }
        \label{tab:real-dsi-esmda}
    }{%
        \begin{tabular}{cccccc}
\toprule
\multirow{2}{*}{Loc.} 
& 
\multicolumn{2}{c}{Observation coverages} & 
\multirow{2}{*}{%
    MSE
} & 
\multicolumn{2}{c}{Data mismatch} 
\\
& History & Validation & & Avg. & Std. dev.
\\
\midrule
No & 0.63 & 0.57 & 2.28 & 1.24 & 0.02 \\
Yes & 0.63 & 0.57 & 2.28 & 1.24 & 0.02 \\
\bottomrule
\end{tabular}

    }{%
        \fonte{Self-elaboration.}
    }
\end{table}

\begin{table}[htbp]
    \centering
    \IBGEtab{%
        \caption{%
            Gated DSI-ESMDA results for the real field in the Brazilian Post-salt with varying gating.
        }
        \label{tab:real-gated-dsi-esmda-1-16}
    }{%
        \begin{tabular}{cccccccccccc}
\toprule
\multirow{3}{*}{
    \begin{minipage}{2cm}
        \centering
        Gating\\
        intensity
    \end{minipage}
}
&
\multirow{3}{*}{Loc.} 
& 
\multicolumn{4}{c}{Observation coverages} & 
\multicolumn{2}{c}{\multirow{2}{*}{
    MSE
}} & 
\multicolumn{4}{c}{Data mismatch} 
\\
& & \multicolumn{2}{c}{History} & \multicolumn{2}{c}{Validation} &     &     & \multicolumn{2}{c}{Avg.} & \multicolumn{2}{c}{Std. Dev.}
\\
& & Pr. & Po.                   & Pr. & Po.                      & Pr. & Po. & Pr. & Po.                & Pr. & Po.
\\
\midrule
\multirow{2}{*}{$\lambda=1$} & No & 0.86 & 0.53 & 0.95 & 0.52 & 2.71 & 2.32 & 5.44 & 1.21 & 2.19 & 0.01 \\
 & Yes & 0.86 & 0.54 & 0.95 & 0.55 & 2.71 & 2.32 & 5.44 & 1.21 & 2.19 & 0.01 \\
\midrule
\multirow{2}{*}{$\lambda=2$} & No & 0.86 & 0.56 & 0.95 & 0.55 & 2.71 & 2.36 & 5.44 & 1.27 & 2.19 & 0.02 \\
 & Yes & 0.86 & 0.57 & 0.95 & 0.58 & 2.71 & 2.36 & 5.44 & 1.27 & 2.19 & 0.02 \\
\midrule
\multirow{2}{*}{$\lambda=4$} & No & 0.86 & 0.63 & 0.95 & 0.64 & 2.71 & 2.44 & 5.44 & 1.52 & 2.19 & 0.04 \\
 & Yes & 0.86 & 0.64 & 0.95 & 0.66 & 2.71 & 2.44 & 5.44 & 1.53 & 2.19 & 0.05 \\
\midrule
\multirow{2}{*}{$\lambda=8$} & No & 0.86 & 0.73 & 0.95 & 0.74 & 2.71 & 2.55 & 5.44 & 2.15 & 2.19 & 0.19 \\
 & Yes & 0.86 & 0.73 & 0.95 & 0.75 & 2.71 & 2.55 & 5.44 & 2.16 & 2.19 & 0.19 \\
\midrule
\multirow{2}{*}{$\lambda=16$} & No & 0.86 & 0.80 & 0.95 & 0.81 & 2.71 & 2.66 & 5.44 & 3.10 & 2.19 & 0.57 \\
 & Yes & 0.86 & 0.80 & 0.95 & 0.81 & 2.71 & 2.66 & 5.44 & 3.12 & 2.19 & 0.57 \\
\bottomrule
\end{tabular}

    }{%
        \fonte{Self-elaboration.}
    }
\end{table}

Mild gating intensities ($\lambda=1$, $\lambda=2$) produce similar results to the standard DSI-ESMDA.
This indicates that the gating mechanism may act as a substitute to simpler localizations on complexely correlated features, 
which is a significant advantage since localization tuning isn't straightforward.
